\section{Design Overview}

Ghostext is easiest to read as a two-layer system. The packet layer turns plaintext into an authenticated binary payload. The token-selection layer embeds that payload into ordinary model generation under a shared prompt. Keeping these layers separate matters because the two layers solve different problems. The packet layer is responsible for confidentiality, integrity, and payload normalization. The token-selection layer is responsible for synchronized replay under one concrete model/tokenizer/runtime configuration.

\paragraph{Packet layer.} Alice first derives internal keys from the shared passphrase and fresh salt, then wraps the plaintext into an opaque packet with authenticated encryption. The released implementation separates internal header/body keys, avoids a plaintext-visible magic marker in the packet bootstrap, and binds decoding to a configuration fingerprint before plaintext release. In practice, this means Bob does not attempt partial or heuristic recovery: if the packet decrypts and authenticates, he gets the exact plaintext; if shared state has drifted or the text was modified in transit, decoding aborts. This packetization step also makes the carried bitstream look pseudorandom given the passphrase, which is the intended input for the embedding layer.

\paragraph{Replay contract.} Ghostext's sender and receiver must share more than a passphrase. They must also share model/tokenizer identity, prompt text, seed, candidate-policy settings, quantization parameters, and codec budgets. The implementation records this contract explicitly through backend metadata and a configuration fingerprint. Operationally, Ghostext treats replay mismatches as protocol errors rather than soft decoding noise. That choice is stricter than ordinary text generation, but it is central to the system claim of deterministic recoverability: Alice and Bob either follow the same path or the protocol stops.

\paragraph{Embedding path.} Once Alice has a packet, she grows the cover text token by token. At each next-token step, Ghostext starts from the same truncated candidate set that natural sampling would see under the chosen \texttt{top\_p} and \texttt{max\_candidates} policy. It then removes candidates that are not retokenization-stable under the receiver's replay path, quantizes the remaining probability mass into integer frequencies, and uses interval narrowing to select one token according to the next payload bits. When the payload is exhausted, the encoder may append a short natural tail so the text can end more naturally. Bob replays the same steps, reconstructs the packet, and ignores any trailing natural tail once packet recovery is complete.

\paragraph{Why the split matters.} Relative to a bare plaintext-to-bits pipeline, packet framing changes both security wording and systems behavior. On the security side, it lets us talk about an authenticated, plaintext-normalized payload rather than about direct leakage from plaintext bit patterns. On the systems side, it introduces explicit packet overhead and makes synchronized replay a first-class requirement. These are not side details: the evaluation in \S8 should be read as evidence about the full two-layer system, not only about arithmetic coding in isolation.
